\Chapter{22}

\Verse{1}
{
	\hE{וַיְהִי אַחַר הַדְּבָרִים הָאֵלֶּה וְהָאֱלֹהִים נִסָּה אֶת אַבְרָהָם וַיֹּאמֶר אֵלָיו אַבְרָהָם וַיֹּאמֶר הִנֵּנִי}
}
{
	\eE{
		And it was after these things,
		and God tested Abraham.
		And He said to him, “Abraham.”
		And he said, “I’m here.”
	}
}
{
	After all that, God tests Abraham.

	He goes ``Abraham!''

	And Abe's like ``I'm right here.''

	I guess that's the first test, and he passed.

	Now for number two.
}

\Verse{2}
{
	\hE{וַיֹּאמֶר קַח נָא אֶת בִּנְךָ אֶת יְחִידְךָ אֲשֶׁר אָהַבְתָּ אֶת יִצְחָק וְלֶךְ לְךָ אֶל אֶרֶץ הַמֹּרִיָּה וְהַעֲלֵהוּ שָׁם לְעֹלָה עַל אַחַד הֶהָרִים אֲשֶׁר אֹמַר אֵלֶיךָ}
}
{
	\eE{
		And He said, “Take your son,
		your only one, whom you love,
		Isaac,
		and go to the land of Mo-riah
		and make him a burnt offering there
		on one of the mountains that I’ll say to you.”
	}
}
{
	So he says ``Take your son''

	(Begin Midrash)

	Abe: Which son?

	(End Midrash)

	God: Your only one.

	(Begin Midrash)

	Abe: I have two sons.

	(End Midrash)

	God: Who you love.

	(Begin Midrash)

	Abe: I love them both.

	God: Christ Abe...

	(End Midrash)

	God: Isaac.

	(Begin Midrash)

	Abe: Ok.

	God: Where was I? Oh right. Take Isaac...

	(End Midrash)

	... and go to the land of YHWH's Choice.

	(Begin Midrash)

	Abe: Which one did Ya choose?

	God: That is the choice Abe. It's called Moriah. It means Ya Choice.

	Abe: My choice?

	God: No no, Yah's Choice. It's My Choice. It means Moriah. I mean Moriah means that.

	Abe: Ok, where in Moriah?

	God: Where was I? Oh right. Go there...

	(End Midrash)

	... to one of the mountains \eRJE{of My Choice}

	(Begin Midrash)

	Abe: Which mountain of your choice?

	God: Dammit Abe stop it, I didn't say ``of My Choice,'' that was just to remind you it's in Moriah. It's called a resumptive repetition.

	Abe: A resumptive what?

	God: Repetition. It's a redactor thing.

	Abe: Why red?

	God: What?

	Abe: Why does the actor have to be red?

	God: Redactor Abe, it's like an Editor.

	Abe: Why not just call it the Editor?

	God: Because E was already taken Abe.

	Abe: He was already taken how?

	God: E, Abe! The letter E. It stands for the Elohist.

	Abe: What's an Elohist?

	God: The author of this section Abe. It means the author calls me God.

	Abe: I'm not following.

	God: See how it's Yellow? That's the Elohist. It's a pun. He's the Yellohist.

	Abe: What's yellow?

	God: The bible Abe, for God's sake.

	Abe: What's a b---

	God: Abe forget that, back to the point. Take your son.

	Abe: Right, Isaac.

	God: Go to the land of Ya Choice.

	Abe: Right, Moriah.

	God: Right.

	Abe: To an unspecified mountain.

	God: No, it's a specific one.

	Abe: You haven't decided yet have you?

	God: Of course I have.

	Abe: Which one?

	God: I'm not telling you yet. It's a specific mountain with a specific name...

	(End Midrash)

	... that I'll say to you

	(Begin Midrash)

	Abe: Ok.

	God: Repeat it all back to me so I know you're following.

	Abe: Take Isaac, go to Moriah, then I wait for you to tell me which mountain to go to, then I go to the mountain, and I think that's as far as we got.

	God: Great. Make sense?

	Abe: Makes sense so far. Then what?

	God: Then

	(End Midrash)

	make him a burnt offering

	(Begin Midrash)

	Abe: What kind of burnt offering should I make for him?

	God: I didn't say for.

	Abe: What?

	God: What?

	Abe: Can you repeat the instructions?

	God: Sure.

	Abe: Take Isaac, go to Moriah, figure out which mountain, go there, and then what

	God: Then

	(End Midrash)

	burn him alive.

	(Begin Midrash)

	Abe: Ok.

	(End Midrash)
}

\Verse{3}
{
	\hE{וַיַּשְׁכֵּם אַבְרָהָם בַּבֹּקֶר וַיַּחֲבֹשׁ אֶת חֲמֹרוֹ וַיִּקַּח אֶת שְׁנֵי נְעָרָיו אִתּוֹ וְאֵת יִצְחָק בְּנוֹ וַיְבַקַּע עֲצֵי עֹלָה וַיָּקם וַיֵּלֶךְ אֶל הַמָּקוֹם אֲשֶׁר אָמַר לוֹ הָאֱלֹהִים}
}
{
	\eE{
		And Abraham got up early in the morning
		and harnessed his ass
		and took his two boys with him
		and Isaac, his son.
		And he cut the wood for the burnt offering,
		and he got up and went to the place
		that God had said to him.
	}
}
{
	% \aB{Who are you by the way?}

	% \aA{In what sense?}

	So Abe gets up early in the morning.

	He grabs his ass and his two boys.

	And Isaac.

	He cuts up the wood.

	Then he gets up and goes to the place.
}

\Verse{4}
{
	\hE{בַּיּוֹם הַשְּׁלִישִׁי וַיִּשָּׂא אַבְרָהָם אֶת עֵינָיו וַיַּרְא אֶת הַמָּקוֹם מֵרָחֹק}
}
{
	\eE{
		On the third day:
		and Abraham raised his eyes
		and saw the place from a distance.
	}
}
{
	% \aB{A name would be nice.}

	% \aA{We can't here.}

	Three days later, Abe looks up.

	He sees the place.
}

\Verse{5}
{
	\hE{וַיֹּאמֶר אַבְרָהָם אֶל נְעָרָיו שְׁבוּ לָכֶם פֹּה עִם הַחֲמוֹר וַאֲנִי וְהַנַּעַר נֵלְכָה עַד כֹּה וְנִשְׁתַּחֲוֶה וְנָשׁוּבָה אֲלֵיכֶם}
}
{
	\eE{
		And Abraham said to his boys,
		“Sit here with the ass;
		and I and the boy:
		we’ll go over there,
		and we’ll bow,
		and we’ll come back to you.”
	}
}
{
	% \aB{An initial at least?}

	% \aA{Most I can do.}

	So Abe's like ``Wait here. We'll be back.''
}

\Verse{6}
{
	\hE{וַיִּקַּח אַבְרָהָם אֶת עֲצֵי הָעֹלָה וַיָּשֶׂם עַל יִצְחָק בְּנוֹ וַיִּקַּח בְּיָדוֹ אֶת הָאֵשׁ וְאֶת הַמַּאֲכֶלֶת וַיֵּלְכוּ שְׁנֵיהֶם יַחְדָּו}
}
{
	\eE{
		And Abraham took the wood for the burnt offering
		and put it on Isaac, his son,
		and took the fire and the knife in his hand.
		And the two of them went together.
	}
}
{
	% \aB{So can you?}

	% \aA{Of course}

	He makes his son carry the wood.

	Abe carries the fire and the knife.
}

\Verse{7}
{
	\hE{וַיֹּאמֶר יִצְחָק אֶל אַבְרָהָם אָבִיו וַיֹּאמֶר אָבִי וַיֹּאמֶר הִנֶּנִּי בְנִי וַיֹּאמֶר הִנֵּה הָאֵשׁ וְהָעֵצִים וְאַיֵּה הַשֶּׂה לְעֹלָה}
}
{
	\eE{
		And Isaac said to Abraham, his father;
		and he said, “My father.”
		And he said, “I’m here, my son.”
		And he said, “Here are the fire and the wood,
		but where is the sheep for the burnt offering?”
	}
}
{
	% \aB{Who are you?}

	% \aA{What's in a name}

	So Isaac's like ``...Dad?''

	Abe says ``I'm here my son.''

	And Isaac's like ``Where's the sheep?''
}

\Verse{8}
{
	\hE{וַיֹּאמֶר אַבְרָהָם אֱלֹהִים יִרְאֶה לּוֹ הַשֶּׂה לְעֹלָה בְּנִי וַיֵּלְכוּ שְׁנֵיהֶם יַחְדָּו}
}
{
	\eE{
		And Abraham said,
		“God will see to the sheep for the burnt offering,
		my son.”
		And the two of them went together.
	}
}
{
	% \aB{An initial I said.}

	% \aA{As in the beginning}

	And Abe goes and I quote...

	``God will provide the lamb for the offering my son.''

	And commas hadn't been invented yet so that sentence is as non-associative in the original text as it is here.

	It's like in Oedipus when Oedipus says:

	``You should have found the murderer the king a mighty king had been destroyed.''

    % I think \chineseE font (cjk-wqy-zenhei.ttc) supports hangul.
	Got it really great I mean the best writing does stuff like that this that is We got it good.\footnote{According to the latest version of The Translator's Transㄴational Transʾ\heb{ל}atiʿn TransLation (To/\heb{ל})(ʾl)kit, the meaning of the above sentence when rendered into the more familiar orthography of Standard Written Academic Typography is: IDENTIFICATION RECURSION: Maximum stack depth exceeded while building parse tree. Please run \texttt{sudo --sudordinal --limit-ordinal ulimit -s }$\aleph_0$ to increase the runtime stack size of child processes to the smallest transfinite cardinal. If you intend to cite this error message in a manuscript to be submitted for publication, please cite the work as:\\IRE of COBAʿAL.©®™ COmmon Biblically \& Akkadʿemically Accurate Language is a registered trademark of BBL's own BBL One! line of translation and transliteration products and server side services. Babel™ runs on Gen 0b10, where \texttt{emerge @world} takes only 0o111 days! Thank you for choosing BBL\chineseE{잉}lɪʃ edition. And remember, when our production quality products produce a scarier error message when compiling your compilation than you've ever seen before and you run into one or two of the exceptional edge case exceptions and panics and runtime \chineseC{二}||s written by our world class members, remember: don't run and don't panic: Keep Calm and Babylon.''}
}

\Verse{9}
{
	\hE{וַיָּבֹאוּ אֶל הַמָּקוֹם אֲשֶׁר אָמַר לוֹ הָאֱלֹהִים וַיִּבֶן שָׁם אַבְרָהָם אֶת הַמִּזְבֵּחַ וַיַּעֲרֹךְ אֶת הָעֵצִים וַיַּעֲקֹד אֶת יִצְחָק בְּנוֹ וַיָּשֶׂם אֹתוֹ עַל הַמִּזְבֵּחַ מִמַּעַל לָעֵצִים}
}
{
	\eE{
		And they came to the place
		that God had said to him.
		And Abraham built the altar there
		and arranged the wood,
		and he bound Isaac, his son,
		and put him on the altar
		on top of the wood.
	}
}
{
	% \aB{Spit it out.}

	% \aA{Not that way}

	So Abe puts down the firewood and ties up his son and he puts his tied up son on the wood like anyone would.
}

\Verse{10}
{
	\hE{וַיִּשְׁלַח אַבְרָהָם אֶת יָדוֹ וַיִּקַּח אֶת הַמַּאֲכֶלֶת לִשְׁחֹט אֶת בְּנוֹ}
}
{
	\eE{
		And Abraham put out his hand
		and took the knife
		to slaughter his son.
	}
}
{
	% \aB{What do they call you?}

	% \aA{Who are you calling they}

	And he takes the knife to slaughter his son.
}

\Verse{11}
{
	\hRJE{וַיִּקְרָא אֵלָיו מַלְאַךְ יְהֹוָה מִן הַשָּׁמַיִם וַיֹּאמֶר אַבְרָהָם אַבְרָהָם וַיֹּאמֶר הִנֵּנִי}
}
{
	\eRJE{
		And an angel of YHWH\eRed{\footnote{Fourth of five. This one's above my pay grade.}} called to him from the skies
		and said, “Abraham! Abraham!”
		And he said, “I’m here.”%
	}
}
{
	\eRJE{
		An' just then, an angel of Jehovah calls to 'im from the 'eavens loik an angel.
		An' 'e says to 'im, “Abraham! Abraham!”
		An' Abraham's, 'e says, “Wot're you on about? I'm roight 'ere, mate!”%
	}%
	\footnote{
		Lit. \eRJE{
			/ən dʒəsʔ ðen ən ˈeɪndʒəl əv dʒəˈhəʊvə kɔːlz tə ɪm frəm ði ˈɛvənz
			lɔɪk ən ˈeɪndʒəl ən i sɛz tə ɪm | ˈeɪbrəˌam ˈeɪbrəˌam | ən ˈeɪbrəˌamz i sɛz |
			wɒʔ ə juː ɒn əˈbaʊʔ | aɪm rɔɪʔ ɪə meɪʔ/
		}
		in the original text, translated here via the Latin alphabet into something
		roughly like English, for simplicity.
	}
}

\Verse{12}
{
	\hRJE{וַיֹּאמֶר אַל תִּשְׁלַח יָדְךָ אֶל הַנַּעַר וְאַל תַּעַשׂ לוֹ מְאוּמָה כִּי עַתָּה יָדַעְתִּי כִּי יְרֵא אֱלֹהִים אַתָּה וְלֹא חָשַׂכְתָּ אֶת בִּנְךָ אֶת יְחִידְךָ מִמֶּנִּי}
}
{
	\eRJE{
		And he said,
		“Don’t put your hand out toward the boy,
		and don’t do anything to him,
		because now I know that you fear L\hspace{-0.22em}Ꞁ,
		and you didn’t withhold your son,
		your only one, from me.”%
	}
}
{

	\eRJE{
		An' 'e says, “Oi! Don't you lay a bloody finger on that lad!
		Don't do nuffin' to 'im! I can see now you fear God, can't I?
		You was actually gonna do it! Your own son an' everyfink!”%
	}%
	\footnote{
		Lit. \eRJE{
			/ən i sɛz | ɔɪ | dəʊnʔʃə leɪ ə ˈblʌdi ˈfɪŋɡər ɒn ðæʔ læd
			| dəʊnʔ duː ˈnʌfɪn tə ɪm | aɪ kən siː naʊ jə fɪə ɡɒd kɑːn aɪ |
			jə wəz ˈækʃəli ˈɡɒnə duː ɪʔ | jər əʊn sʌn ən ˈɛvrɪfɪŋk/
		}
		in the original text.
	}
}

\Verse{13}
{
	\hRJE{וַיִּשָּׂא אַבְרָהָם אֶת עֵינָיו וַיַּרְא וְהִנֵּה אַיִל אַחַר נֶאֱחַז בַּסְּבַךְ בְּקַרְנָיו וַיֵּלֶךְ אַבְרָהָם וַיִּקַּח אֶת הָאַיִל וַיַּעֲלֵהוּ לְעֹלָה תַּחַת בְּנוֹ}
}
{
	\eRJE{
		And Abraham raised his eyes and saw,
		and here was a ram behind,
		caught in the thicket by its horns.
		And Abraham went and took the ram
		and made it a burnt offering
		instead of his son.%
	}
}
{

	% \aB{You?}
	% \aA{And that would make you}

	\eRJE{
		So Abraham looks up an' 'ere, wot's this then?
		There's a ram caught by 'is 'orns in a bush!
		“Well, that'll do nicely,” says Abraham, an' 'e
		goes over, grabs the ram, an' offers it up instead
		of 'is boy.%
	}%
	\footnote{
		Lit. \eRJE{/səʊ ˈeɪbrəˌam lʊks ʌp ən ɪə | wɒʔs ðɪs ðen
		| ðɛəz ə ræm kɔːʔ baɪ ɪz ɔːnz ɪn ə bʊʃ | wɛl ðæʔəl duː
		ˈnaɪsli | sɛz ˈeɪbrəˌam | ən i ɡəʊz ˈəʊvə ɡræbz ðə ræm
		ən ˈɒfəz ɪʔ ʌp ɪnˈstɛd əv ɪz bɔɪ/} in the original text.
	}

}

\Verse{14}
{
	\hRJE{וַיִּקְרָא אַבְרָהָם שֵׁם הַמָּקוֹם הַהוּא יְהֹוָה יִרְאֶה אֲשֶׁר יֵאָמֵר הַיּוֹם בְּהַר יְהֹוָה יֵרָאֶה}
}
{
	\eRJE{
		And Abraham called the name of that place
		“YHWH Yir’eh,”
		as is said today:
		“In YHWH’s mountain it will be seen.”%
	}
}
{

	\eRJE{
		An' Abraham calls the place “Jehovah'll Sort It,”
		which is why folks still say, “Up on Jehovah's mountain,
		'e'll sort it.”%
	}%
	\footnote{
		Lit. \eRJE{ən ˈeɪbrəˌam kɔːlz ðə pleɪs dʒəˈhəʊvəl sɔːʔ ɪʔ
		| wɪtʃ ɪz waɪ fəʊks stɪl seɪ | ʌp ɒn dʒəˈhəʊvəz ˈmaʊnʔɪn
		| iːl sɔːʔ ɪʔ/} in the original text.
	}

}

\Verse{15}
{
	\hRJE{וַיִּקְרָא מַלְאַךְ יְהֹוָה אֶל אַבְרָהָם שֵׁנִית מִן הַשָּׁמָיִם}
}
{
	\eRJE{
		And an angel of YHWH called to Abraham
		a second time from the skies%
	}
}
{

	\eRJE{
		Then the angel of Jehovah calls down to Abraham from the 'eavens again.
	}%
	\footnote{
		Lit. \eRJE{/ðen ði ˈeɪndʒəl əv dʒəˈhəʊvə kɔːlz daʊn tə ˈeɪbrəˌam frəm ði ˈɛvənz əˈɡɛn/}
		in the original text.
	}
}

\Verse{16}
{
	\hE{וַיֹּאמֶר בִּי נִשְׁבַּעְתִּי}
	\hRJE{נְאֻם יְהֹוָה}
	\hE{כִּי יַעַן אֲשֶׁר עָשִׂיתָ אֶת הַדָּבָר הַזֶּה וְלֹא חָשַׂכְתָּ אֶת בִּנְךָ אֶת יְחִידֶךָ}
}
{
	\eE{
		and said, “I swear by me
	}%
	\eRJE{
		—word of YHWH—
	}%
	\eE{
		that because you did this thing
		and didn’t withhold your son,
		your only one,
	}
}
{

	and God's like:
	
	I swear to you Abe%
	\eRJE{—An' these 'ere are the words o' Jehovah—}%
	\footnote{Lit. \eRJE{/ən ðiːz ɪə ɑː ðə wɜːdz ə dʒəˈhəʊvə/} in the original.}%
	that because you did this thing, because you didn't withhold your son, your only one,
}

\Verse{17}
{
	\hE{כִּי בָרֵךְ אֲבָרֶכְךָ וְהַרְבָּה אַרְבֶּה אֶת זַרְעֲךָ כְּכוֹכְבֵי הַשָּׁמַיִם וְכַחוֹל אֲשֶׁר עַל שְׂפַת הַיָּם וְיִרַשׁ זַרְעֲךָ אֵת שַׁעַר אֹיְבָיו}
}
{
	\eE{
		that I’ll bless you
		and multiply your seed
		like the stars of the skies
		and like the sand that’s on the seashore,
		and your seed will possess its enemies’ gate.
	}
}
{

	that I'll spread your {\cuma} all over the place Abe, for real this time, I really will.

	I'll spread your {\cuma} so much it'll break open the door to your enemies town and then your {\cuma} will rush inside like a rushing river of {\cuma} and drown all your enemies Abe I promise.
}

\Verse{18}
{
	\hE{וְהִתְבָּרְכוּ בְזַרְעֲךָ כֹּל גּוֹיֵי הָאָרֶץ עֵקֶב אֲשֶׁר שָׁמַעְתָּ בְּקֹלִי}
}
{
	\eE{
		And all the nations of the earth
		will be blessed through your seed
		because you listened to my voice
	}
}
{

	And all the nations of the earth are gonna get blessed by your seed just like that Abe.

	Because you listened.

	You really listened Abe.

	God bless you.

}

\Verse{19}
{
	\hE{וַיָּשׁב אַבְרָהָם אֶל נְעָרָיו וַיָּקֻמוּ וַיֵּלְכוּ יַחְדָּו אֶל בְּאֵר שָׁבַע וַיֵּשֶׁב אַבְרָהָם בִּבְאֵר שָׁבַע}
}
{
	\eE{
		And Abraham went back to his boys
		And they got up
		And went together to Beer-sheba
		And Abraham lived in Beer-sheba
	}
}
{
	Then Abraham \aB{And Isaac?} gets up \aB{What about Isaac?} and goes back to his boys.

	And they went together to Seven Wells.

	And Abe lives in Seven Wells.
}

\Verse{20}
{
	\hJ{וַיְהִי אַחֲרֵי הַדְּבָרִים הָאֵלֶּה וַיֻּגַּד לְאַבְרָהָם לֵאמֹר הִנֵּה יָלְדָה מִלְכָּה גַם הִוא בָּנִים לְנָחוֹר אָחִיךָ}
}
{
	\eJ{
		And it was after these things,
		and it was told to Abraham, saying,
		“Here Milcah has given birth, she also,
		to sons for your brother Nahor:
	}
}
{
	After all that, someone tells Abe:

	Hey, Milcah had the baby.

	You're an uncle Abe.
}

\Verse{21}
{
	\hJ{אֶת עוּץ בְּכֹרוֹ וְאֶת בּוּז אָחִיו וְאֶת קְמוּאֵל אֲבִי אֲרָם}
}
{
	\eJ{
		Uz, his firstborn,
		and Buz his brother,
		and Kemuel, the father of Aram,
	}
}
{
	\Table{|p{0.18\linewidth}|p{0.16\linewidth}|p{0.56\linewidth}|}{
		\hline
		\textbf{Name}
			& \textbf{Hebrew}
			& \textbf{Meaning / likely etymology} \\
		\hline
		\textbf{Uz}
			& \heb{עוּץ}
			& Uncertain. Possibly related to \heb{עוץ}, \textbf{``counsel / advise.''} \\
		\hline
		\textbf{Buz}
			& \heb{בּוּז}
			& \textbf{``Contempt / despised.''} From \heb{בוז}, ``despise.'' \\
		\hline
		\textbf{Kemuel}
			& \heb{קְמוּאֵל}
			& Probably \textbf{``El/God has risen''} or ``raised by El,'' from \heb{קום},
				 ``rise.'' \\
		\hline
		\textbf{Aram}
			& \heb{אֲרָם}
			& \textbf{Aram / Aramea.} Historical etymology uncertain; traditionally associated
				 with \heb{רום}, \textbf{``high / exalted.''} \\
		\hline
	}
}

\Verse{22}
{
	\hJ{וְאֶת כֶּשֶׂד וְאֶת חֲזוֹ וְאֶת פִּלְדָּשׁ וְאֶת יִדְלָף וְאֵת בְּתוּאֵל}
}
{
	\eJ{
		and Chesed and Hazo
		and Pildash and Jidlaph
		and Bethuel.
	}
}
{
	\Table{|p{0.18\linewidth}|p{0.16\linewidth}|p{0.56\linewidth}|}{
		\hline
		\textbf{Name}
			& \textbf{Hebrew}
			& \textbf{Meaning / likely etymology} \\
		\hline
		\textbf{Chesed}
			& \heb{כֶּשֶׂד}
			& Probably the eponym of the \textbf{Chaldeans} (\emph{Kasdim}); original meaning
				 uncertain. \\
		\hline
		\textbf{Hazo}
			& \heb{חֲזוֹ}
			& Probably related to \heb{חזה}, \textbf{``see / behold,''} perhaps ``vision.'' \\
		\hline
		\textbf{Pildash}
			& \heb{פִּלְדָּשׁ}
			& \textbf{Unknown.} Probably a non-Hebrew personal name. \\
		\hline
		\textbf{Jidlaph}
			& \heb{יִדְלָף}
			& \textbf{``He drips / weeps.''} From \heb{דלף}, ``drip.'' \\
		\hline
		\textbf{Bethuel}
			& \heb{בְּתוּאֵל}
			& Probably \textbf{``dwelling/house of El''} or similar; exact etymology disputed. \\
		\hline
	}
}

\Verse{23}
{
	\hJ{וּבְתוּאֵל יָלַד אֶת רִבְקָה שְׁמֹנָה אֵלֶּה יָלְדָה מִלְכָּה לְנָחוֹר אֲחִי אַבְרָהָם}
}
{
	\eJ{
		And Bethuel fathered Rebekah.”
		Milcah gave birth to these eight
		for Nahor, Abraham’s brother.
	}
}
{
	\Table{|p{0.18\linewidth}|p{0.16\linewidth}|p{0.56\linewidth}|}{
		\hline
		\textbf{Name}
			& \textbf{Hebrew}
			& \textbf{Meaning / likely etymology} \\
		\hline
		\textbf{Bethuel}
			& \heb{בְּתוּאֵל}
			& Probably \textbf{``dwelling/house of El''} or similar; exact etymology disputed. \\
		\hline
		\textbf{Rebekah}
			& \heb{רִבְקָה}
			& \textbf{Uncertain.} Often connected with a root involving \textbf{binding / tying},
				 perhaps ``tether.'' \\
		\hline
		\textbf{Milcah}
			& \heb{מִלְכָּה}
			& \textbf{``Queen.''} Feminine of \heb{מֶלֶךְ}, ``king.'' \\
		\hline
		\textbf{Nahor}
			& \heb{נָחוֹר}
			& \textbf{Uncertain.} Probably an old Semitic personal name; sometimes connected with
				 a root meaning ``snort/breathe heavily,'' but insecure. \\
		\hline
		\textbf{Abraham}
			& \heb{אַבְרָהָם}
			& Traditionally \textbf{``father of a multitude,''} from the biblical wordplay
				 \emph{av hamon} (\heb{אַב הֲמוֹן}); exact historical etymology uncertain. \\
		\hline
	}
}

\Verse{24}
{
	\hJ{וּפִילַגְשׁוֹ וּשְׁמָהּ רְאוּמָה וַתֵּלֶד גַּם הִוא אֶת טֶבַח וְאֶת גַּחַם וְאֶת תַּחַשׁ וְאֶת מַעֲכָה}
}
{
	\eJ{
		And his concubine,
		whose name was Reumah,
		she too gave birth,
		to Tebah and Gaham
		and Tahash and Maacah.
	}
}
{
	\Table{|p{0.18\linewidth}|p{0.16\linewidth}|p{0.56\linewidth}|}{
		\hline
		\textbf{Name}
			& \textbf{Hebrew}
			& \textbf{Meaning / likely etymology} \\
		\hline
		\textbf{Reumah}
			& \heb{רְאוּמָה}
			& Probably \textbf{``exalted / raised up,''} from \heb{רום}, ``be high.'' \\
		\hline
		\textbf{Tebah}
			& \heb{טֶבַח}
			& \textbf{``Slaughter / butchery.''} From \heb{טבח}, ``slaughter.'' \\
		\hline
		\textbf{Gaham}
			& \heb{גַּחַם}
			& \textbf{Unknown / uncertain.} \\
		\hline
		\textbf{Tahash}
			& \heb{תַּחַשׁ}
			& \textbf{``Tahash.''} Also the biblical word for an uncertain animal or type of hide
				 used for the Tabernacle; deeper etymology uncertain. \\
		\hline
		\textbf{Maacah}
			& \heb{מַעֲכָה}
			& Probably \textbf{``crushed / pressed,''} from \heb{מעך}, ``crush, press.'' \\
		\hline
	}
}
